%% kapa: exam=1 round=25 year=2014 period=1 form=A subject=economics subject_ko=경제원론 seq=34 number=74
%% source: Q-Net 감정평가사 기출문제 (cst00309) artlSeq=5101207 file=extracted/1008/5101207/2014년 제25회 감정평가사 1차시험 기출문제/1교시(A형).hwp
%% conversion: kordoc 4.24.51 -> tools/gen1.py (auto); review=proofread-sonnet; flags=none
\begin{question}{74}
총공급곡선이 $Y = \overline{ Y } + \alpha (P-P^{ e } )$인 총수요-총공급 모형에서 경제가 현재 장기균형상태에 있다. 이 경제의 중앙은행이 통화량을 감소시킬 경우, 물가 예상이 합리적으로 형성되고 통화량 감소가 미리 예측된다면 다음 설명 중 옳은 것은? (단, $Y$는 실질GDP, $\overline{ Y }$는 실질GDP의 장기균형 수준, $\alpha$는 0보다 큰 상수, $P$는 물가, $P^{ e }$는 예상물가수준이다.)
\begin{choices}
  \item 실질GDP는 즉시 감소한 다음 서서히 원래 수준으로 복귀한다.
  \item 물가는 즉시 감소한 다음 서서히 원래 수준으로 복귀한다.
  \item 물가는 즉시 감소하고 실질GDP도 즉시 감소한다.
  \item 물가는 서서히 감소하고 실질GDP는 즉시 감소한다.
  \item 물가는 즉시 감소하고 실질GDP는 원래 수준을 유지한다.
\end{choices}
\end{question}
\answer{5}
